\section{张量网络元数的集中}\label{sec:local}

我们的目标是在保持张量网络收缩结果的同时，把其中的大表集中到少数几个站点。这里的预算是各站点元数的平方和。例如，\(W/T\) 个元数为 \(T\) 的站点，其预算为 \(WT\)。若每次把 \(D\) 个站点合并成元数为 \(DT\) 的站点，预算就会增至原来的 \(D\) 倍。引理~\ref{lem:local} 将大站点的预算控制在 \(O(WT)\)，同时把其个数降至 \(O(W/(DT))\)。于是它们的总元数为 \(O(W/\sqrt D)\)。这一变换还会引入小表。第四部分将在相继两次集中操作之间缩短小表网络的连线，再通过批量查表计算所得网络的值。

一个\emph{站点}是由有序比特列表上的任意一张表给出的张量。其元数既计入外部比特，也计入站点所接内部连线的端点。把站点分组，分别置于一个最大度为 \(g\) 的宿主图的各个单元中。连线可以位于一个单元内部、连接相邻单元，或通向外部。同一个宿主图可以承载多个网络。若两个网络都使用同一单元，则将它计数两次。被指定为\emph{不可移动}的站点保留其外部比特，并留在原单元中。所有外部比特都必须属于这样的站点。

\begin{lemma}[局部集中]\label{lem:local} 固定一个满足 \(\omega_k<9/4\) 的有限域 \(k\)、一个度上界 \(g\) 和一个整数 \(\ell\ge1\)。对每个充分大且固定的整数 \(D\)，以下结论对一切充分大的整数 \(T\) 成立。取域扩张 \(K/k\)，使它包含某个固定的有限扩张，且 \(|K|\) 至少为 \(T\) 的某个充分大的多项式，并满足 \([K:k]=O_D(\log T)\)。输入表可以已经定义在 \(K\) 上。

设被占用的单元总数至多为 \(W/(DT)\)，不可移动站点的总数至多为 \(W/(DT)\)。在每个网络中，每个单元至多有 \(100D\) 个站点，总元数至多为 \(100DT\)，且对所有网络的全部站点有 \(\sum_i t_i^2\le WT\)。每个网络的张量都可以在 \(K\) 上表示成若干具有相同外露比特的网络之和。从每个网络的和式中任取一项，所选各项合在一起满足以下界。\begin{enumerate}
\item 共有 \(O(W/(DT))\) 个\emph{冻结站点}，其中包括在原单元内持有原外部比特的站点。其最大元数为 \(O(DT)\)，元数平方和为 \(O(WT)\)，总元数为 \(O(W/\sqrt D)\)。\item 其余站点合称\emph{核心}，各自的元数至多为 \(D^{-\ell}T\)，总元数为 \(O(W(1+\log D)^2)\)。\item 在同一个新宿主图中，用 \(D^{O(1)}\) 个位置替换每个原单元，可能的连接只发生在同一单元内或相邻单元间。核心位置诱导的宿主子图的度有绝对常数上界。任意两个冻结站点之间均无直接连线。存在 \(c=c(g,\ell,k)\)，使所有启用的内部线束的宽度均可取为至少 \(\lceil T/D^c\rceil\)。即使对每个站点均至少按该元数计费，前两项结论中的界仍然成立。\item 外层求和在每个被占用的单元使用 \(O_D(\log(T+1))\) 个控制比特。为给出计算界，用标签表示每张输入表，用有序线束表示其端口。每对站点对应一束，用长度和两端的槽位描述，同一束两端的比特顺序一致。任何局部重排都是给定表标签的一部分。假设给定一个使用 \(O_D(1)\) 种颜色的固定着色，使距离在某个仅依赖 \(D,g\) 的界以内的不同单元颜色各异。在旧宿主图上经过 \(O_D(1)\) 轮通信，即可判定指标是否合法、确定输出布局并追踪外露子列表。当每个输入标签不超过 \(A\) 比特时，每个被占用单元的工作量和通信量分别为 \(O_D(A+(1+\log T)^{O(1)})\) 次比特运算和同一数量级的比特。每张输出表在一个由输入字典与参数确定的字典中有长度为 \(O_D(A+\log(T+1))\) 的标签。\end{enumerate}
前三项结论中的隐含常数可依赖于 \(k,g,\ell\)，但与 \(D,T,W\) 无关。若将输入界中的数值 \(100\) 改成另一个固定数值，则只有这些常数随之改变。\end{lemma}

证明在每个被占用的单元中使用两个球面形式。第一个形式通过成对的张量腿连接相邻单元，每一对称为一条\emph{接缝}。它用标记点处的一个插入来表示一个站点。该插入的能量控制编码其状态所需宽度的平方。通过让插入穿过接缝并在共同的标记点处合并，可将其能量集中到冻结站点。接缝另有自己的能量预算。每次穿越均耗去其中一部分，最后将每条接缝分配给许多小表。

开始移动前，必须先实现输入网络实际给出的连线。第二个球面始终留在本单元内，用于实现这些连线。根据推论~\ref{cor:universal}，它的插入可以编码第一个球面上指定的连线方式。两个球面的腿都参与合并，但只有第一个球面的腿可以穿过接缝。两个球面使用不同的物种数和能量尺度，从而使穿越后留在原处的私有腿保持较小的宽度。

\subsection{球面形式与一个转置恒等式}\label{subsec:seam}

我们采用第~\ref{sec:fermions} 节的代数自由费米子构造，这里使用多个物种。用 \(V^{(r)}_e\) 表示如下外代数中能量为 \(e\) 的子空间。其生成元为 \(\psi_{s,m}\)，物种为 \(1\le s\le r\)，模态为 \(m\ge0\)，能量为 \(m+\tfrac12\)。按 \((s,m)\) 的字典序排列生成元，取相应的楔积基。对于每个基向量，其长度都同余于 \(2e\pmod2\)。

令 \(c_0=\pi/\sqrt3\)。当 \(q>0\) 时，单物种的生成函数为 \(Z(q)=\prod_{m\ge0}(1+\exp(-q(m+1/2)))\)。由于 \(t\mapsto\ln(1+\exp(-t))\) 是凸函数，逐一在长度为 \(q\) 的相邻区间上积分，得到 \(\ln Z(q)\le q^{-1}\int_0^\infty\ln(1+\exp(-t))\,dt=\pi^2/(12q)\)。对于 \(r\) 个物种，有 \(\dim V^{(r)}_e\le\exp(qe)Z(q)^r\)。当 \(e>0\) 时对 \(q\) 取最小值，当 \(e=0\) 时使用真空态，便得
\begin{equation}\label{eq:local-width}
      \dim V^{(r)}_e\le \exp(c_0\sqrt{re}).
\end{equation}

我们用基矢的二进制编号为每条腿编码，并将未用编码的权重设为零。

给定互异点 \(z_i\in K\)，球面形式 \(C_r\) 将一组楔积映为一个普法夫式。按位点的次序列出所有生成元。同一位点内的配对以及不同物种间的配对都等于零。对同一物种，在互异位点 \(i,j\) 依此次序取模态 \(m,n\)，二者的配对为 \[
 [x^m y^n]\,(z_j+y-z_i-x)^{-1}.
\] 空矩阵的普法夫式为一，奇数长度的列表给出零。交换能量分别为 \(e,d\) 的两个相邻齐次位点会引入符号 \((-1)^{(2e)(2d)}\)。因此，当能量固定时，这些符号都只是固定的标量。将整系数交替配对公式约化以后，这些定义也适用于特征二。

设依次位于 \(z,z+u\) 的相邻两位点的能量分别为 \(e,m\)。将二者替换为位于 \(z\)、能量为 \(e'\) 的一个位点，再预合成下列映射，就得到 \(u^{e'-e-m}\) 的系数：\[
 I_{e'}:V^{(r)}_e\otimes V^{(r)}_m\longrightarrow V^{(r)}_{e'}.
\] 只有满足 \(e'-e-m\in\mathbb Z\) 的 \(e'\ge0\) 才会出现。计算 \(I_{e'}\) 时，对偏移量分别为 \(0,1\) 的两位点之间的部分内部配对求和，将未匹配的生成元在零点展开以进行平移，最后在能量 \(e'\) 处依原顺序取楔积。更准确地说，偏移量为 \(a\) 的平移将 \(\psi_{s,d}\) 替换为 \(\sum_{f\ge d}\binom fd a^{f-d}\psi_{s,f}\)。这就是命题~\ref{prop:fusion} 的融合公式在多物种情形下的推广。其证明把完整匹配分解成内部配对与外部配对，可对每个物种分别进行。将此形式看作任一位置的函数时，其极点只能位于其他位置，且在无穷远处有界。

要连接两个球面，就在各指定的一个位点上，以能量 \(e\) 的楔积基上的相等配对收缩两条腿。我们称这对腿为一条\emph{接缝}。要让插入穿过接缝，我们需要相对于这种普通的相等配对对融合系数作转置。

\begin{lemma}[接缝处的转置]\label{lem:seam} 设 \(a,b\) 分别为能量 \(e,e'\) 的基矢，并设 \(x\in V_m^{(r)}\)。对 \(j\le m\) 存在保持奇偶性的映射 \(S_j:V_m^{(r)}\to V_j^{(r)}\)，它们与 \(a,b\) 无关；且存在非零标量 \(\alpha=\alpha(e,e',m)\)，使得
\begin{equation}\label{eq:seam-adjunction}
 [I_{e'}(a,x)]_b
    =\alpha\sum_{j\le m}[I_e(b,S_jx)]_a .
\end{equation} 系数均相对于所示楔积基取出。\end{lemma}
\begin{proof} 在 \(k\) 的固定扩张中加入 \(-1\) 的一个平方根 \(i\)。取 \(y\) 为无穷远处的局部坐标，因而对应位置为 \(1/y\)。在取 \(y\) 的系数以提取模态之前，该处每出现一个生成元，就将核乘以一次 \(i/y\)。于是有限坐标 \(U\) 与此坐标图之间的配对为 \(i/(1-Uy)\)。因此，零点与无穷远点处的二点形式对长度为 \(d\) 的相同基矢给出配对系数 \[
      (-1)^{d(d-1)/2}i^d=
\begin{cases}1& d\ \text{ 为偶数},\\ i& d\ \text{ 为奇数},\end{cases}
\] 其余系数全为零。考虑在 \(0,1,\infty\) 三点取值于 \(a,x,b\) 的三点形式，首先取零点与一点之间的内部匹配。将剩余的生成元与无穷远处的生成元配对，恰好实现定义 \(I_{e'}(a,x)\) 的平移和楔积。因此，三点形式的值等于它的 \(b\) 分量乘以上式中的奇偶性因子。

用 \(V=1/U\) 对有限坐标取逆，并在每个变元上乘因子 \(-i/U\)，二点核便保持不变。事实上，\((-i/U_1)(-i/U_2)/(1/U_2-1/U_1)=1/(U_2-U_1)\)。原来的无穷远坐标图变成一个因子为一的普通有限坐标图，原来的零点则变成无穷远坐标图，每个生成元各贡献一个负号。在中间位点，取如下单生成元代换的外幂：\[
 \psi_{s,d}\longmapsto
 \sum_{f=0}^{d}
 \left([X^d]\,\frac{-i}{1+X}
       \left(\frac{1}{1+X}-1\right)^{\!f}\right)\psi_{s,f}.
\] 投影至能量 \(j\) 即得 \(S_j\)。每一项均保持长度，且能量至多为 \(m\)。倒转位点顺序，令 \(b\) 位于零点、\(a\) 位于无穷远点，再计算三点形式。倒序引入的符号、来自原零点的符号以及两个对角奇偶性因子均只由原先三个奇偶性决定。取相应因子之比即证得~\eqref{eq:seam-adjunction}。\end{proof}

以单物种为例，取 \(a=1\) 和 \(b=x=\psi_0\)。在能量 \(1/2\) 处，近侧融合将 \(x\) 原样传过。取逆后，中间插入为 \(-i\psi_0\)；在远侧，取偏移量分别为 \(0,1\) 的两个位点与 \(b\) 作相应收缩，在零能量处给出 \(-i\)。乘以因子 \(\alpha=i\)，就恢复了原来的系数一。同样交换输入与输出的能级，便可使一般插入穿越接缝。

\section{在标记点布设连线}\label{subsec:planting}

输入网络中的一个站点可能与许多其他站点有连线。为了将该站点作为单个插入来移动，我们把这些连线在该站点的各端点编码到一条腿中。此时，球面形式必须能检验每条连线两端的比特是否一致。我们分别在每个网络的每个被占用单元中实现这一检验。

每个单元有两个球面形式。形式 $A$ 承载站点的腿以及通向相邻单元的接缝。$A$ 上的额外标记点提供一个检验连线的联合态。私有形式 $B$ 通过每个额外标记点上单独的一张插入表实现该联合态。因此，布设完毕后，每个插入在每个球面上恰有一条腿，此外还可带有外露比特或控制比特。在原站点和接缝处，$B$ 腿均处于真空态。图~\ref{fig:stack} 展示了这两类插入。

\begin{figure}[htb]\centering
\begin{tikzpicture}[x=1cm,y=1cm,font=\small,
  strip/.style={draw=ink,rounded corners=3pt,line width=.65pt},
  ins/.style={draw=ink,rounded corners=2pt,fill=white,inner xsep=7pt,inner ysep=4pt},
  leg/.style={draw=ink,line width=.8pt}]
\draw[strip,fill=ink!10] (-4.35,.98) rectangle (4.35,1.55);
\draw[strip,fill=ink!5] (-4.35,-1.55) rectangle (4.35,-.98);
\node[anchor=west] at (-4.22,1.26) {缝合球面 \(A\)};
\node[anchor=east] at (4.22,1.26) {\(R\) 个物种};
\node[anchor=west] at (-4.22,-1.26) {私有球面 \(B\)};
\node[anchor=east] at (4.22,-1.26) {\(R/b\) 个物种};
\draw[leg] (-5.35,1.26)--(-4.35,1.26);
\draw[leg] (4.35,1.26)--(5.35,1.26);
\node[above,font=\footnotesize] at (-4.85,1.26) {接缝};
\node[above,font=\footnotesize] at (4.85,1.26) {接缝};
\node[ins] (sta) at (-2.15,0) {原站点};
\node[ins] (aux) at (2.0,0) {辅助插入};
\draw[leg] (sta.north)--(-2.15,.98);
\draw[leg,densely dotted] (sta.south)--(-2.15,-.98);
\node[anchor=west,font=\footnotesize] at (-2.1,-.6) {真空};
\draw[leg] (aux.north)--(2.0,.98);
\draw[leg] (aux.south)--(2.0,-.98);
\node[anchor=west] at (2.05,.62) {\(z^P\)};
\node[anchor=west] at (2.05,-.62) {\(z\)};
\foreach \x/\y in {-2.15/.98,2.0/.98,2.0/-.98}
    \fill[ink] (\x,\y) circle (1.5pt);
\draw[draw=ink,fill=white] (-2.15,-.98) circle (1.5pt);
\end{tikzpicture} \caption{站点表把与该站点相接的内部连线的比特编码在 \(A\) 上。辅助表将 \(A\) 与 \(B\) 耦合，共同施加各连线两端相等的约束。只有 \(A\) 与相邻单元缝合。}\label{fig:stack}
\end{figure}

$B$ 上位于 $z$ 的标记点对应 $A$ 上位于 $z^P$ 的标记点，其中 $P$ 是域的特征的一个幂。将 $z$ 移动 $t$，就使 $z^P$ 移动 $t^P$。因此，以 $A$ 上的洛朗阶为单位时，这一对腿的能量为 $e_A+P^{-1}e_B$。给定同一能量预算，$B$ 上允许的能量是前者的 $P$ 倍。我们在 $B$ 上使用较少的物种，以使对应的腿保持较小的宽度。这将给出站点穿越接缝时留在原处的腿的宽度上界。

按引理~\ref{lem:local} 固定 $g,\ell$。令 $F$ 为不小于 $D^{1000(\ell+1)}$ 的最小特征幂，且令 $a=\lceil F^3\rceil$。给定组数 $M$ 及充分大的特征幂 $b>F$，取 $R$ 为不超过 $T/a$ 的最大 $bM$ 倍数。令
\begin{equation}\label{eq:plant-scales}
 R'=R/b,\qquad P=b/F,\qquad \kappa=P^{-1},\qquad H=Ra^2 .
\end{equation} 其中 $R$ 和 $R'$ 分别是 $A$ 和 $B$ 上的物种数。当 $T$ 充分大时，$T/(2a)\le R\le T/a$。若 $e_A+\kappa e_B\le h$，则由式~\eqref{eq:local-width}，两侧宽度的上界分别为 $\sqrt{Rh}$ 和 $\sqrt{R'Ph}=\sqrt{Rh/F}$ 的常数倍。以比特为单位，忽略取整误差，便有
\begin{equation}\label{eq:plant-widths}
 w_A=O\bigl(T\sqrt{h/H}\bigr),\qquad
 w_B=O\bigl(T\sqrt{h/(FH)}\bigr).
\end{equation}

\subsection{单个物种的连线形式}

将每对不同站点之间的连线合为一束有序线束，把端点处的置换并入站点表。先在这些表内收缩自环。把每条内部线束拆分成每包 $R$ 比特的数据包，末包不足处补零。物种 $s$ 将承载每个包中的第 $s$ 个比特。每个站点与 $O_g(D)$ 束线相接，每个单元连同与之关联的边总共涉及 $O_g(D^2)$ 束线。因此，在元数为 $t_i$ 的站点处，单个物种至多承载 $t_i/R+O_g(D)$ 比特。

选取整数能量 $h_i$ 处不同的楔积基矢来编码站点 $i$ 处的比特串。根据引理~\ref{lem:dimension}，沿整数能级趋于无穷时，单物种空间维数的对数渐近等于 $(\pi/\sqrt3)\sqrt{h_i}$。由于 $R\ge T/(2a)$ 且 $a\ge D$，可以取 \[
                 h_i=O_g\bigl(a^2(1+t_i/T)^2\bigr).
\] 在接缝 $j$ 处，跨过交界面的各包对一个物种共贡献 $O_g(DT/R+D^2)=O_g(Da)$ 比特，故整数能量 $q_j=O_g(D^2a^2)$ 已经足够。两个单元在此处使用相同的比特次序与基矢编码。固定 $D$ 时，每个物种的包数、能量与编码均取自有限集。

在这些站点与接缝的能级空间上，用系数定义多线性形式 $L$。当所有编码均有效且每一比特的两端均一致时，赋予相应基元组系数一。对于通向另一单元的连线，将它的站点端与接缝端相比较。其余系数均为零。我们将在固定的\emph{总}能量处缝合多物种接缝。为恢复每个物种各自正确的能量，令 $Q=\sum_j q_j$，并将 $L$ 在其余满足总和至多为 $Q$ 的非负半整数接缝能量元组上延拓为零。各站点能量保持固定。

\begin{lemma}[权重有界的辅助标记点]\label{lem:plant-probes} 在单物种球面上多项式个额外标记点插入一个联合态，就能实现延拓后的形式 $L$。在每个额外标记点，此态只使用真空与 $\psi_0$。其支撑集内每个码字所占用的标记点数均为偶数且至多为 $O_g(Da)$。\end{lemma}
\begin{proof} 长度为 $d$ 的楔积，其能量至少为 $\tfrac12+\tfrac32+\cdots+(d-\tfrac12)=d^2/2$。各站点处的总长度为 $O_g(a\sum_i(1+t_i/T))=O_g(Da)$。对于至多 $g$ 条接缝，由柯西–施瓦茨不等式及其总能量界 $Q=O_g(D^2a^2)$，也得出相同的长度上界。列出所得输入空间中出现的原标记点与模态的一切数对 $(v,m)$。此列表的大小至多为 $D,a$ 的多项式。

在新点 $z$ 处插入 $\psi_0$ 时，它与上述模态的配对由坐标为 $(z-z_v)^{-m-1}$ 的余向量给出。在任意特征下，这些不同的主部给出线性无关的有理函数。选取互异的正则点，使其余向量张成模态空间的对偶空间。事实上，若某个非零线性组合在所有已选点上均为零，它必在代数闭包的某个未使用的点取非零值。所需点数至多等于列表中的模态数，且这些点都在某一有限扩张内。如有必要可添入真空，使辅助标记点数至少为三。

对由 $d$ 个新标记点构成的一个子集，在这些点上插入 $\psi_0$，在其余点上插入真空。先考虑所得普法夫式中每个辅助生成元均与一个原生成元配对的部分。该部分本身就是至多 $d$ 个标记点上的插入检验的线性组合。事实上，把所有辅助生成元排在末尾，再按每个匹配所含辅助生成元内部配对的确切集合进行分组。每个非空配对集合的贡献等于其收缩所得的标量，乘以一个仅依赖于配对位置的符号，再乘以剩余标记点上没有辅助内部配对的部分。归纳地减去这些贡献，就分离出所需部分。

现取输入空间直和的偶部中的一个非零向量，令 $d$ 为其非零分量的最大楔积总长度。用 $d$ 个标记点且不含辅助内部配对的检验在所有更短的楔积上均为零。在长度为 $d$ 的分量上，它所计算的是选定 $d$ 个余向量的外积在该分量上的值，仅相差一个共同的符号。各个原基元组在以数对 $(v,m)$ 的列表为生成元的外代数中，是互不相同的基楔积。因此，总有一种 $d$ 个余向量的选法能检出该分量。由此得到的偶检验张成偶部的整个对偶空间。实际使用的每个插入都具有不超过 $d=O_g(Da)$ 的偶权重，并在奇向量上取零。规定的整数能量给出偶长度，因此 $L$ 的零延拓在奇部上也取零。插入检验的一个线性组合遂给出所需联合态。\end{proof}

在 $R$ 个物种中各用一份该态。在站点 $i$ 处，插入表将每个物种限制在能级 $h_i$，对数据包解码，检查填充比特为零，并给出带有外露比特的原表项。在接缝 $j$ 处，我们对总能量 $Rq_j$ 处相同的基矢实行缝合。为看出这能恢复原来的连线，将普法夫式按物种分组。先将原标记点和辅助标记点排为前后相接的两个列表，再将每个列表的元素按物种重排。辅助列表重排所引入的符号可并入联合态。由于每种物种的辅助权重均为偶数，将两个列表交错排列不会产生额外符号。原列表重排的符号只依赖输入楔积。在收缩辅助态后所得的任何非零乘积中，由零延拓的定义，每一物种的接缝能量元组或为指定元组，或其和大于 $Q$。所有物种的这些和加起来等于 $RQ$，所以每一物种都必须取指定元组。此时全部原楔积长度均为偶数，剩余的重排符号为一。该乘积恰好检验了每条连线两端的比特是否相等。我们还在每个单元内放置一个真空汇集点，供移动步骤使用。

\subsection{实现辅助联合态}\label{subsec:plant-private}

引理~\ref{lem:plant-probes} 指定了所有辅助标记点共享的一个张量。为了移动，我们需要在每个标记点上分别放置一张插入表。私有球面通过局部映射给出这一共享张量。一个 $B$ 物种将对应 $A$ 上含 $b$ 个物种的一块。

将缝合侧的 $R$ 个物种依次划分成 $R'$ 个连续块。若某块在辅助标记点 $i$ 处有 $d_i$ 个 $\psi_0$ 模态被占用，则称该块的\emph{类型}为 $\tau=(d_i)_{i=1}^m$。将块张量投影到这些局部子空间，得到张量 $J_\tau$，其中计入了块内重排的符号。它的第 $i$ 条腿的维数为 $\binom b{d_i}$。固定由各块类型构成的有序字后，联合张量等于各 $J_\tau$ 的乘积，仅相差一个已知符号。交换来自不同块的两个位点列表，对符号的指数仅贡献两列表长度的乘积模二。只保留满足 $J_\tau\ne0$ 的类型。

用 $H_2$ 表示以比特计的二元熵，其中 $H_2(0)=H_2(1)=0$。取间距为 $1/m$ 的固定正网格，将每个 $H_2(d_i/b)$ 向上舍入到严格大于它的下一个网格点 $u_i$。于是
\begin{equation}\label{eq:plant-ent}
 \log_2\binom b{d_i}\le b u_i,\qquad
 \max_i u_i\le 2,\qquad \sum_i u_i=O_{g,\ell}(Da\log D).
\end{equation} 事实上，由引理~\ref{lem:plant-probes} 有 $s=\sum_i d_i/b=O_g(Da)$。由凹性可知，各熵之和不超过 $s\log_2(em/s)$，并约定在 $s=0$ 时此式取零。固定 $\ell$ 及域的特征时，标记点个数 $m$ 至多为 $D$ 的多项式，且舍入使总和至多增加一。

在对应的 $B$ 标记点上，取所有满足第 $i$ 个位点处能量至多为 $C(bu_i)^2$ 的非零单物种球面形式的直和 $S_\tau$，其中 $C$ 是稍后选取的绝对常数。当 $b$ 充分大时，存在常数 $c_\tau$，使得对每个 $d\ge1$，一组局部线性映射的乘积给出
\begin{equation}\label{eq:plant-comparison}
 \bigoplus_{s=1}^{c_\tau} S_\tau^{\otimes d}
                    \longrightarrow J_\tau^{\otimes d}.
\end{equation} 因为 $c_\tau$ 与 $d$ 无关，同一选择器指标就能供所有出现的该类型共用。

为进行比较，取自然对数，并在代数闭包上工作，那里指数仍严格小于 $9/4$。对位于 $x_1,\ldots,x_m$ 的归一化谱价格，令 $M_u=\inf_o\sum_i u_i\,d(o,x_i)$。在 $o$ 处制备 $J_\tau$，再输运其各条腿，由式~\eqref{eq:plant-ent} 可知，其价格至多为 $b(\log 2)\sum_i u_i d(o,x_i)$。取下确界得到 $b(\log 2)M_u$。对大小为二的对角张量应用相同的输运上界，该张量的价格为 $\log 2$，故得 $\inf_o\sum_i d(o,x_i)\ge1$。因此 $M_u\ge\min_i u_i>0$。取尺度 $N=b$、宽度上限 $c_0\sqrt C\,u_i$，应用推论~\ref{cor:universal}。其中的比较张量恰好是 $S_\tau$，其价格为 $b(c_0\sqrt C\,M_u+o(1))$，且此估计关于归一化价格一致。取 $c_0\sqrt C>2\log 2$。固定 $D$ 时，网格向量和 $A$ 上的位置选择均只有有限种，且均定义在某个固定的有限扩张上。取 $P$ 次根在该域上给出置换，因此当 $b$ 改变时，$B$ 上可能的位置也都落在有限集中。于是，当特征幂 $b$ 充分大时，对每一种价格，$S_\tau$ 都严格支配 $J_\tau$。由第~\ref{sec:prices-background} 节关于域扩张和严格比较的结果，在固定的定义域上得到式~\eqref{eq:plant-comparison}。这些映射可以延拓至输入表所在的域。

固定类型字。对于每种重数为 $d>0$ 的类型，选取一个副本指标 $s$，并为 $d$ 个因子分别选取一个 $B$ 能量元组，依此展开式~\eqref{eq:plant-comparison} 中的限制。按原来的块顺序将各因子分配给私有物种。它们的乘积就是多物种球面形式，仅相差一个置换齐次列表产生的符号。此时，各位点映射给出一张局部表，将该标记点处的 $A$ 腿和 $B$ 腿耦合起来。它可以联合地作用于\emph{该标记点处}同一类型的所有块。对所有选择和类型字求和，便得到所需的辅助张量。原站点和接缝处的 $B$ 腿仍保持真空态。

\paragraph{控制与存储。}一个\emph{私有选择器}记录有序类型字、每一块的 $B$ 能量以及每种出现的类型所对应的一个指标 $s$。用单张表检查这些数据的合法性，包括未用槽位是否补零，并给出置换符号。由于 $d_i\le b$ 且能量的两倍为 $O(b^2)$，每个块记录类型和两倍能量共需 $O(m\log(2b))$ 比特。全部副本指标总计耗费 $O_{D,b}(1)$ 比特。共有 $R/b$ 个块，对任意固定 $\epsilon>0$，先取充分大的 $b$，再取充分大的 $T$，就可使选择器的长度小于 $\epsilon T$。

在 $A$ 上，每个被占用的辅助模态耗费能量 $1/2$，故辅助总能量为 $O_g(RDa)$。在一个 $B$ 物种上，由式~\eqref{eq:plant-ent}，各能量上限的和为 $O(b^2\sum_i u_i^2)=O_{g,\ell}(b^2 Da\log D)$。乘以 $R'$ 并换算为 $A$ 的单位，得到每个单元中的能量界 \[
 \sum_i(e_{A,i}+\kappa e_{B,i})
    =O_{g,\ell}(RDa+F RDa\log D)\le D H
\] 这里取 $D$ 充分大，并用到了 $H=Ra^2$ 与 $a\ge F^3$。就存储而言，使用所选块的实际空间。在标记点 $i$ 处，由式~\eqref{eq:plant-ent} 和~\eqref{eq:local-width}，任一条腿均只需要 $O(bu_i+1)$ 比特。对所有块求和，每个辅助插入的这两条腿共需 $O(R)=O(T/a)$ 比特。再对各标记点求和，总共需 $O_{g,\ell}(DT\log D)$ 比特，其中取 $b\ge m$ 以吸收舍入开销。

选择器在网络内部求和，但外层布局必须知道每个标记点处的能量与宽度。将所有块依序分成 $M$ 个等大的连续组。在外层求和中，为每组指定其各数对（类型，$B$ 能量元组）的直方图。直方图记录各数对的重数，私有选择器记录它们的次序。字母表固定，与 $T$ 无关，因此直方图在每个单元内共需 $O_{D,b,M}(\log(T+1))$ 比特。它们决定每组中各条腿的能量和编码连接后的总长度。表将这些直方图与选择器相互核验，并按照选择器给出的块顺序进行解码。每个合法选择器恰好出现于一个外层求和项中。因此，可以先为分组布局选择 $M$，再为选择器布线选择 $\epsilon$，最后为价格比较及选择器宽度选择 $b$，所有选择都在 $T$ 增长之前完成。

\section{在缝合曲面上移动插入}\label{subsec:motion}

布设构造将原来各表表示在每个单元中的一对球面上。现在我们把这些表合并到少数标记点处，只需利用洛朗展开的两类结果。在插入标记点处，取一个系数，可将移动中的表与原来就在该处的表融合。在接缝处，取一个阶数足够小的负阶系数，可将移动中的表移至相邻单元，并降低接缝能量。我们利用这一下降量，既能限制穿越次数，也能限制任一单元所能接收的能量。

在每个单元中，将原不可移动站点的标记点以及另一个真空标记点指定为\emph{汇集点}。它们留在各自的单元内。依次处理原来的可移动站点，跳过已经被吸收的站点。移动中的表保留其普通比特端口。若该表现有的两条腿在缝合球面 $A$ 和私有球面 $B$ 上的能量分别为 $m_A,m_B$，令 \[
       m=m_A+m_B/P,\qquad \mu=m+H,\qquad L=\lceil 2P\mu\rceil .
\] 其中 $P$ 是式~\eqref{eq:plant-scales} 的换算率。从一次移动开始直至结束，储备能量 $H$ 始终随这次移动同行，而 $m$ 则在每次融合或穿越后更新。

\subsection{由碰撞重建取值}

固定外层求和的一个分支，为推导下述恒等式也暂时固定私有选择器。分支本身就决定了能量，因此同一恒等式适用于所有容许的选择器取值。固定所有其他位置和插入表。收缩所得的每一表项都是移动者的 $B$ 坐标 $z$ 的有理函数，其 $A$ 坐标为 $z^P$。球面核只可能在当前球面上的其他标记点处有极点，且在无穷远处有界。在 $B$ 坐标为 $z_0$ 的标记点附近，令 $s=z-z_0$。由于 $P$ 是特征的一个幂，$A$ 上相应的坐标差为 $s^P$。因此，阶数为 $i$ 的 $A$ 融合与阶数为 $v$ 的 $B$ 融合对收缩的 $j=Pi+v$ 阶项作出贡献。在加权能量为 $m_0$ 的插入处，合并所得的能量为
\begin{equation}
            m_{\rm new}=m_0+m+j/P.                     \label{eq:motion-fusion}
\end{equation} 事实上，由命题~\ref{prop:fusion}，每次融合都将其阶数加到参与融合的两个能量之和上。两个输出能量非负，故 $i$ 与 $v$ 均有下界。因此极点阶数有界，对每个 $j$，只有有限多个数对 $(i,v)$ 作出贡献。

为重建当前位置的取值，使用单元中其他标记点处的下列系数。\begin{itemize}
\item 对每个插入，取其整个主部。\item 在指定的真空汇集点，还取 $0,\ldots,gL$ 阶各项，即使该汇集点已经吸收了一张表，仍然这样做。\item 在每条接缝处，限于所给的极点阶数界，仅取 $j\le -L$ 的各项。\end{itemize}
若函数的所有这些系数都为零，则它的极点只可能出现在至多 $g$ 条接缝处，且每处的极点阶数至多为 $L-1$。它在指定汇集点的零点阶数至少为 $gL+1$，在无穷远处没有极点。射影直线上非零有理函数的除子次数为零，故此函数恒为零。因此，在具有规定极点界的函数空间上，于移动者当前位置取值的泛函是上述系数泛函的线性组合。此组合逐表项作用，并且与各张表的具体内容无关。

插入处的一个系数是对所有相容数对 $(i,v)$ 的求和。对每一数对，施行两个融合映射，并将它们与移动表及静止表收缩，计入所选次序的符号。与汇集点碰撞便结束本次移动。在任何其他插入处只使用 $j<0$，因此合并后的表继续移动时所带的加权能量至多为原先两者之和。这一步可以吸收一个辅助插入、一个尚未轮到的原站点，或一个由先前穿越所留下的私有插入。

\subsection{穿越接缝}

设在能量为 $e$ 的一条接缝处选取了 $j\le-L$ 阶的系数。只有球面 $A$ 在此处缝合。接缝处 $B$ 上的插入是真空，因此 $B$ 展开的泰勒阶数为 $v\ge0$。对这样的每个 $v$，若 $i=(j-v)/P$ 为整数且落在融合范围内，则接缝处新的 $A$ 能量满足
\begin{equation}
       e'=e+m_A+i\le e+m_A-2\mu\le e-\mu .             \label{eq:motion-drop}
\end{equation} 满足 $e'<0$ 的项均为零。固定 $i$，用此前未使用的正则位置上的取值对 $B$ 的泰勒系数进行插值。在所得每一项中，近侧留下一个能量为 $m_B$ 的私有插入，其在 $A$ 上为真空。一条普通的相等连线将它的指标传至原表。我们称这一私有插入为本次穿越的\emph{残留插入}。

为了确定新的移动表，由线性性可以固定其原来的 $A$ 态 $x$。令 $a,b$ 分别遍历能量 $e,e'$ 的楔积基矢，并以 $C_{\rm near},C_{\rm far}$ 表示余下的两个球面形式，其中略去未参与本步变换的变元。引理~\ref{lem:seam} 给出恒等式 \[
\begin{aligned}
 &\sum_b C_{\rm near}(b)\sum_a [I_{e'}(a,x)]_b\,C_{\rm far}(a)\\
 &\hspace{1em}=\alpha\sum_{t\le m_A}\sum_b C_{\rm near}(b)
                      \sum_a [I_e(b,S_t x)]_a\,C_{\rm far}(a).
\end{aligned}
\] 在右端，内层求和是在远侧 $A$ 球面上插入 $S_t x$ 趋近 $b$ 时的 $e-e'-t$ 阶洛朗系数。该阶数为整数，因为非零项中的 $e'-e-m_A=i$ 与 $m_A-t$ 都是整数。用远侧正则位置上的取值对此系数进行插值，在相应的 $B$ 位置取真空。对 $b$ 求和就在能量 $e'$ 处将两球面缝合。新的移动表现位于远侧，能量为 $(t,0)$；它保留普通端口以及通向残留插入的连线（图~\ref{fig:motion-crossing}）。所有重排符号都只依赖指定的能量，即使 $x$ 与私有指标或某个未参与变换的变元相耦合，恒等式依然成立。

\begin{figure}[!htbp] \centering
\begin{tikzpicture}[x=1cm,y=1cm, font=\small,
 cell/.style={draw=ink!50,fill=ink!3,rounded corners=2pt},
 tab/.style={draw=ink,fill=white,rounded corners=1pt,inner sep=3.2pt,
             font=\small,align=center}]
 \node[font=\small\bfseries] at (-3.62,1.45) {近侧单元};
 \node[font=\small\bfseries] at (3.62,1.45) {远侧单元};
 \draw[cell] (-5.36,-1.11) rectangle (-1.88,1.13);
 \draw[cell] ( 1.88,-1.11) rectangle ( 5.36,1.13);
 \node[align=center] at (-3.62,.55) {缝合球面 $A$};
 \node[align=center] at ( 3.62,.55) {缝合球面 $A$};
 \fill[sea] (-1.88,.55) circle (2pt);
 \fill[sea] ( 1.88,.55) circle (2pt);
 \draw[sea,line width=.9pt] (-1.88,.55) -- (1.88,.55);
 \node[fill=white,inner sep=2pt,text=sea] at (0,.55) {接缝能量 $e'$};
 \node[tab] (res) at (-3.62,-.40) {私有残留插入\\ $(0,m_B)$};
 \node[tab] (mov) at ( 3.62,-.40) {移动表\\ $(t,0)$};
 \draw[dashed,rust,line width=.75pt] (res.east)--(mov.west);
 \node[fill=white,inner sep=2pt,text=rust] at (0,-.40) {$B$ 态指标};
\end{tikzpicture} \caption{能量为 $(m_A,m_B)$ 的移动插入从近侧单元穿越到远侧单元。实线表示将能量为 $e'\le e-\mu$ 的接缝状态彼此等同。虚线把原来的私有态传给新的移动表，该表在缝合球面上的能量为 $t\le m_A$。} \label{fig:motion-crossing}
\end{figure}

两次插值均可在原收缩所在的同一个有限域 $K$ 上完成。乘以一个公分母，就能将极点受指定限制且在无穷远处有界的有理函数空间嵌入次数不超过总极点阶数上界的多项式空间。在多于此数的正则域点上取值给出单射，所以这些取值泛函张成对偶空间，并能恢复各个所需的洛朗系数。排除所有已被占用的位置。映射 $z\mapsto z^P$ 是 $K$ 上的一个置换，因此可以在两个球面上使用对应位置作正则取值。下面的估计将次数控制在 $O_D(T)$，将每个单元内占用的位置数控制在 $D$ 的某个多项式以内。我们可以先在代数闭包上取得这些界，再据以在进行任何插值前选定 $K$。固定域元素的一种排列顺序。对依该顺序排列的每组容许正则点，在表字典中存入某个重建矩阵的各表项。局部控制只需指定所选点并检查它是否可用。

每次移动至多允许 $U=D^5$ 次穿越。达到此限时，移动者当前的标记点成为汇集点。在相邻两次穿越之间，每次非终止的融合都会消去一个静止插入。一次移动结束后不会再放出新的移动者，因此对每个求和项，该过程都将终止。

\subsection{移动的能量与连线估计}

先估计最终各插入存储的能量。元数为 $t_i$ 的原站点的初始加权能量为 $O((1+t_i/T)^2H)$。辅助插入在每个被占用的单元至多贡献 $DH$。另外，给每个原可移动站点分配一份储备 $H$。在每个单元中，初始插入能量和储备能量之和为 $O(D^2H)$。计入该批次的全部网络，对其所有单元求和所得的界为 \[
 O\!\left(H\left(D\,\#\{\text{被占用的单元}\}
                  +\#\{i\}+T^{-2}\sum_i t_i^2\right)\right)
       =O(WH/T).
\]

由式~\eqref{eq:motion-fusion}，可继续移动的融合不会增加参与融合的能量之和。只保留当前移动的储备。在一次穿越中，能量 $m$ 足以支付远侧的 $t\le m_A$ 和近侧残留插入的 $m_B/P$。在终止移动的融合中，汇集点可能获得的能量至多为 $m+gL/P\le(1+3g)\mu$，因为 $H\ge1$ 且 $P\ge1$。对汇集点的能量赋予权重 $1/(1+3g)$，对其他能量（包括储备）赋予权重一。在上述三种操作中以及移动达到穿越次数上限时，加权总和均不会增加。因此在所得每个求和项中，插入的总能量为 $O_g(WH/T)$。

每个单元内也有相应的界。一次穿越带入的能量计费至多为 $t+H\le\mu$，而由式~\eqref{eq:motion-drop}，接缝损失的能量至少有这么多。一个单元至多有 $g$ 条接缝，每条接缝的初始能量为 $O_g(D^2H)$。再加上该单元的初始计费量，就得到每一步中该单元当前插入能量的上界 $O_g(D^2H)$。此外，每次穿越都使接缝能量至少减少 $H$，因此在整个构造中，与一个单元关联的穿越共发生 $O_g(D^2)$ 次。只有这些穿越可以引入新的静止残留插入。由于 $H=O_D(T)$ 且 $P$ 固定，这些能量界与标记点个数界就给出了所述的总极点阶数上界 $O_D(T)$。在所有单元中，达到次数上限而结束的移动次数至多为穿越总次数除以 $U$。计入最初的不可移动站点与指定的汇集点，汇集点总数为 $O(W/(DT))$。

还须布设展开产生的普通连线。由式~\eqref{eq:plant-widths}，每次穿越留下的私有态占 $O_g(DT/\sqrt F+1)$ 比特。一端跟随移动者，另一端停留在残留插入处，直至被另一插入吸收，再随其移动。仍可能被吸收的残留插入只携带其自身那次穿越所产生的端点。移动端点一旦完成一次移动，其汇集点便固定不动。重复访问也计数时，一次移动至多访问 $U+1$ 个单元，每次访问至多收集 $O_g(D^2)$ 个残留端点，且整次移动至多产生 $U$ 个自身的端点。因此在任意一步跳移中，它所携带的端点数至多为 $O_g((U+1)D^2)$。利用与单元相关联的穿越次数界，并将终止情况计入，可得稍宽的界，每个单元有 $O_g((U+1)D^6)$ 段局部连线。单一单元内连续的移动只需一条局部连接。对第~\ref{subsec:planting} 节所取的 $F$，这些线段在该单元的总宽度为 $O(D^{-2\ell-1}T)+O_D(1)$。

私有选择器的连线从读取选择器的标记点，以及校验选择器和计算球面形式的站点出发。其端点沿相同的路径移动。每个单元初始拥有的此类端口数至多为 $D$ 和组数的某个多项式，因此由路径长度与穿越次数界可知，与一个单元关联的选择器线段及其副本的数量也受同一组参数的多项式控制。该计数与 $b$ 和选择器的编码率无关。固定组数后，选择足够低的编码率，使每个单元中的选择器总宽度（包括执行相等复制所需的宽度）至多为 $D^{-2\ell}T$。下一节将把剩余能级换算为各表的元数。


\section{新元数的估计}\label{subsec:local-output}

移动结束后，每个汇集点都有两条齐次球面腿，也可能还持有外部比特或细连线。接缝仍可能具有较高能量。我们把物种分组，使每组都只接一张小表，即使某一组获得某条接缝的全部能量，仍是如此。

选取 \(M\) 为 \(D\) 的某个幂，具体大小如下指定。在式~\eqref{eq:plant-scales} 中，可以使两个物种数 \(R,R'\) 都被 \(M\) 整除。把每个球面的物种依序分为 \(M\) 个等大的连续组，私有球面的构造中每个完整的块均归入同一组。按各组的能量元组，把每条剩余齐次腿展开成直和。在每个未经改动的辅助插入处使用指定元组，在接缝两侧使用相同元组。每个基矢在展开式中恰出现一次。固定各组的能量后，球面形式分解为各组形式的乘积。重排生成元只引入一个符号，由每个位点处各组的奇偶性确定。

考虑一条一般的腿，它有 \(r\) 个物种、能量为 \(e\)。用 \(e_j\) 表示第 \(j\) 组的能量，则 \(\sum_j e_j=e\)。由式~\eqref{eq:local-width} 与柯西–施瓦茨不等式，总编码宽度满足
\begin{equation}\label{eq:local-chunk-size}
 \sum_{j=1}^M\left\lceil\frac{c_0}{\ln 2}\sqrt{(r/M)e_j}\right\rceil
 \le \frac{c_0}{\ln 2}\sqrt{r e}+M .
\end{equation} 布局中也可以使用整数宽度 \(\lceil4\sqrt{(r/M)e_j}\rceil\)。由于 \(2e_j\) 是整数，这只需计算取整的整数平方根；上述界仍然成立，只需把系数改为四。

未经改动的辅助腿使用其选择器所指定的实际乘积子空间。在一个组中，可由各块类型及其能量的计数读出总比特数，因为对类型为 \((d_i)_i\) 且取整后熵界为 \(u_i\) 的块，第 \(i\) 条腿至多需要 \(O(bu_i)\) 比特。固定选择器后，这些编码单射到相应齐次空间的张量积中。组形式的表对其解码，并对未用的码字赋予零权重。

对每个单元中的每一组，在第一个球面和私有球面上分别用一张任意表实现相应的形式。每个单元的腿数至多为 \(D\) 的多项式。将式~\eqref{eq:local-chunk-size} 用于单独一组，该组的一般端口宽度至多为 \(O(DT/\sqrt M+1)\)；实际辅助子空间的界则表明，每条辅助腿在每组的宽度为 \(O(R/M+1)\)。因此，在加入细的选择器连线前，其中一张表的总元数至多为 \(D^{c_1}T/\sqrt M+O_D(1)\)。指数 \(c_1\) 的选择可独立于 \(M,b\) 及选择器的编码率。将 \(M\) 取为 \(D\) 的一个足够高次幂，使上述量小于 \(D^{-\ell}T/10\)。对这一 \(M\)，取足够低的选择器编码率，使之满足第~\ref{subsec:motion} 节的布线界，然后在布设时选取块大小 \(b\)。依此次序，在 \(T\) 增长之前即可完成全部三个选择。

\paragraph{冻结站点。}各汇集点成为冻结站点。在加权能量为 \(e\) 的一个汇集点处，由式~\eqref{eq:plant-widths} 与~\eqref{eq:local-chunk-size} 可知，两条球面腿的总宽度为 \(O(T\sqrt{e/H})+O(M)\)。移动的界表明，在全部汇集点上有 \(\sum e=O(WH/T)\)，而在每个汇集点处有 \(e=O(D^2H)\)。因此，宽度中的能量项平方和为 \(O(WT)\)，最大值为 \(O(DT)\)。由于各不可移动站点都保留自己的外部端口，各汇集点外部总宽度的平方和由输入界控制。在每个单元中，移动与选择器引入的所有额外连线总宽度为 \(O(D^{-2\ell}T)\)。即使这些连线全部终止于同一汇集点，其在该单元的汇集点上增加的宽度平方代价也至多为 \(O(D^{-4\ell}T^2)\)。固定 \(D\) 时，组数为常数；只要 \(T\) 充分大，取整所引入的加性代价也服从相同的界。用 \((x+y+z)^2\le3(x^2+y^2+z^2)\) 合并这三类代价，即得到包含原外部宽度在内的元数平方和界。由汇集点个数与柯西–施瓦茨不等式，得到 \[
 \sum_{\text{冻结站点}}t_i\le
 \left(O\!\left(\frac W{DT}\right)\,O(WT)\right)^{1/2}
 =O(W/\sqrt D).
\]

\paragraph{核心站点。}为估计各组表的总元数，先将其端口宽度对所有组求和。由第~\ref{subsec:plant-private} 节，未经改动的辅助插入在每个单元中的总编码宽度为 \(O(DT(1+\log D))\)。每个单元至多有 \(g\) 个接缝端口，在式~\eqref{eq:local-chunk-size} 中对各组求和后，总宽度为 \(O_g(DT)\)。每个单元另有 \(O_g(D^2)\) 条残留的私有腿。根据式~\eqref{eq:plant-widths} 和~\eqref{eq:local-chunk-size}，每条腿在全部组中共耗费 \(O_g(DT/\sqrt F)+O(M)\)，所以 \(T\) 充分大时，每个单元的全部残留腿共耗费 \(O(T)\)。与汇集点配对的腿所贡献的元数至多为全局冻结元数。移动时产生的普通连线、选择器连线及所有标量表，在每个单元中共贡献 \(O(T)\)。因此，在控制宿主图的度之前，核心总元数为 \(O(W(1+\log D))\)。未经改动的各辅助插入的元数为 \(O(T/a)\)，各残留插入的元数为 \(O(DT/\sqrt F)\)；各组表和布线宽度均已比 \(D^{-\ell}T\) 小一个固定余量。这些表连同选择器校验表构成全部核心表。

为构造共同的宿主图，在每个单元内给能够在同一求和项中同时出现的各种表和线束预留槽位。每个单元需要 \(D^{O(1)}\) 个槽位。块大小会改变表项与选择器长度，但所有块类型共用同一组槽位。将标量因子在各自单元内合并。起初，允许同一单元或相邻单元里的槽位任意相连。若汇集点直接相连，则在沿其路径携带的细线上插入一个核心恒等表，将直接连接分开。给每个核心站点附上一棵固定的平衡二叉树，每种可能线束对应一个叶子。树根接收该表的有序端口，每个内部顶点将其比特列表分成分别送往两个子节点的子列表。相应的相等性检验表将每个比特沿长度为 \(O(1+\log D)\) 的唯一路径传送，所有内部顶点的元数均至多为原站点元数的两倍。这就证明了所述的度界和核心总元数界。最后，使 \(c\) 超出这些计数中的各个指数，超出的余量取决于 \(\ell\)。于是，对非空线束填充至宽度 \(\lceil T/D^c\rceil\)，或者对每个站点至少按这一元数计费，每个单元至多增加 \(O(D^{-2\ell}T)\) 比特。唯一的补零方式保持张量不变。在上述计数中，它也足以容纳每个单元内有界个仅提供权重的表。这证明了引理~\ref{lem:local} 的前三项结论。

\paragraph{外层求和的指标。}用端点的槽位、长度以及在两端有序端口中的位置编码每一束线。输入标签已经包含全部局部比特置换。沿线束传递有序列表时，无须记录各个外部比特的地址。每个被占用单元的分支信息包括各组最初的直方图、移动与融合的选择（位置、能量、插值点和阶数）以及最终的分组拆分。固定 \(D\) 后，其条目数有界，直方图字母表的大小也有界。每个整数的取值范围为 \(T^{O_D(1)}\)。域中一点用 \(O_D(\log T)\) 比特表示。所有私有选择器比特都是求和项内部的连线，不计入这一数目。将未使用的条目置零，便得到项指标的唯一填充。

使用给定的宿主图的 \((4D^5+20)\) 次幂的着色。固定 \(g\) 时，贪心着色可用 \(O_D(1)\) 种颜色做到这一点。按颜色以及站点槽位的次序依次启动移动。每次移动至多穿越 \(D^5\) 次，故并行的移动影响互不相交的邻域。因此，补齐后的完整调度长度仅由 \(D,g\) 控制。每一步中，各单元均可将自己的局部描述发送给邻居。合并时，将两张输入表的描述与操作参数连接起来。穿越时，添上所指定的插值、融合以及其细连线的新端点。所得任一表的历史只涉及有界条这样的记录，界取决于 \(D\)。因此，一个长 \(O_D(A+\log(T+1))\) 的标签，就能通过固定字典把每个表项指定为端口的函数。该字典可以采用任意插值权重、限制映射与内部收缩。生成标签时无须计算表项的值。

还须验证，上述短计算始终不必检查选择器比特。各组直方图中的计数确定了能量与编码宽度。范围及总和校验容许每个可实现的直方图；若某个通过校验的直方图无法实现，选择器校验将给出零。插值时可以使用由所有正则域点组成的固定有序集合，所用权重则以局部极点界为指标。调度器只需检查是否避开标记点、指标是否在范围内，以及融合与移动是否遵循相应规则。系数为零的项也可保留。各项检查与长度估算所需的只是对 \(O_D(\log(T+1))\) 比特整数或域元素坐标作算术运算，包括带余除法和平方根取整。调度器将旧标签视为字符串，复制到固定的填充槽位。每个单元内记录的数目和对记录做的操作数均有与 \(T\) 无关的上界，从而比特运算次数与通信比特数均为 \(O_D(A+(1+\log T)^{O(1)})\)。这就证明了引理~\ref{lem:local} 的最后一项结论。
